\label{ch:spectral}

\section{Operatorial representation of Dirichlet characters}

The constants in this chapter are grouped because they share a spectral
construction independent of their decimal expansions and free of target
values. Let

\[
Ne_n=ne_n,\qquad n\in\N,
\]

in \(\ell^2(\N)\). A residual character \(\chi\) defines the diagonal operator \(Q_\chi e_n=\chi(n)e_n\); whenever the power is trace class,

\begin{equation}
\Tr(N^{-s}Q_\chi)=L(s,\chi).
\label{eq:L-trace}
\end{equation}

The equation preserves the modulus, orientation, and provenance of each
character; therefore, it does not identify distinct residual objects.

\section{HMT construction of the residue characters}

The preceding diagonal operator is not the initial algebraic datum. The
selection chain is

\[
\begin{aligned}
\APP&\longrightarrow\TRIT\longrightarrow\TPK
\longrightarrow\{\text{residue modulo }3,\ C_P\}\\
&\longrightarrow\{\chi_{-3},\ J=C_P\bmod3\}
\longrightarrow\{\chi_{-3},\ J^2=-I,\chi_{-4}\}\\
&\longrightarrow Q_3,Q_4.
\end{aligned}
\]

The triadic residual functional determines the difference between the two nonzero
oriented classes. The incidence sector uses the symmetric Paley conference
matrix of order six,
\[
C_P^{\mathsf T}=C_P,\qquad C_PC_P^{\mathsf T}=5I_6.
\]
Its reduction \(J=C_P\bmod3\) satisfies
\[
J^2=2I_6=-I_6
\quad\text{in }\operatorname{End}_{\mathbb F_3}(\mathbb F_3^6).
\]
The torsional quarter then determines the cyclic phase
\(1,\ii,-1,-\ii\). Applying the Chinese remainder theorem preserves
both residual data and produces \(Q_{12}=Q_3Q_4\). The trace
\(\Tr(N^{-s}Q)\) is evaluated at the terminal stage; neither the Dirichlet series
nor its decimal expression enters the selection of the character.

The construction is made explicit by the tables
\begin{equation}
\chi_{-3}(n)=
\begin{cases}
0,&3\mid n,\\
1,&n\equiv1\pmod3,\\
-1,&n\equiv2\pmod3,
\end{cases}
\qquad
\chi_{-4}(n)=
\begin{cases}
0,&2\mid n,\\
1,&n\equiv1\pmod4,\\
-1,&n\equiv3\pmod4.
\end{cases}
\label{eq:characters34-tables}
\end{equation}
En \(\ell^2(\N)\),
\[
Q_3e_n=\chi_{-3}(n)e_n,
\qquad Q_4e_n=\chi_{-4}(n)e_n.
\]
As they are diagonals, they commute; the Chinese remainder theorem gives
\begin{equation}
(Q_3Q_4)e_n=\chi_{-3}(n)\chi_{-4}(n)e_n
=\chi_{12}(n)e_n.
\label{eq:q12-product-proof}
\end{equation}
The product simultaneously preserves the triadic residue and the quarter of
torsion. Consequently, the character modulo twelve is determined
before performing any evaluation of its values \(L\).

\begin{remark}[Negative control]
The chain is invalidated if \(Q_3\) does not distinguish the two oriented classes,
if \(J^2\ne-I\), if \(Q_3Q_4\) does not have the declared CRT table, or if the
operator is selected from the value of the constant.
\end{remark}

\section{Internal complex structure and Catalan's constant}

The internal complex structure satisfies

\begin{equation}
J^2=-I.
\label{eq:J-square}
\end{equation}

The equality \(J^2=-I\) defines a representation
\(\rho:C_4\to\operatorname{Aut}_{\mathbb F_3}(\mathbb F_3^6)\),
\(\rho(j)=J\). To realize it spectrally, take the unitary character
\(\upsilon:C_4\to U(1)\), \(\upsilon(j)=\ii\), and the phase map
\(n\mapsto j^n\). On \(\ell^2(\N)\), the induced diagonal representation
is
\[
U_4e_n=\upsilon(j^n)e_n=\ii^ne_n.
\]
Its oriented part is
\[
Q_4=\frac{U_4-U_4^*}{2\ii},\qquad
Q_4e_n=\sin\!\left(\frac{\pi n}{2}\right)e_n
=\chi_{-4}(n)e_n.
\]
This is the explicit lift between the finite generator of order four and the
residual operator; \(\mathbb F_9\) is not identified with \(\mathbb C\).

On odd integers, the same character can be written

\[
\chi_J(n)=\ii^{n-1}=\chi_{-4}(n).
\]

Equivalently,

\begin{equation}
Q_4=\frac{U_4-U_4^*}{2\ii},
\qquad Q_4e_n=\chi_{-4}(n)e_n.
\label{eq:Q4}
\end{equation}

The Catalan constant appears as

\begin{equation}
G=\Tr(N^{-2}Q_4)=L(2,\chi_{-4})
=\sum_{n\ge0}\frac{(-1)^n}{(2n+1)^2}.
\label{eq:catalan}
\end{equation}

The unitary lift of the torsion quarter generates exactly the
character whose quadratic moment is Catalan's constant.

\begin{proof}[Operational proof]
The operator \(N^{-2}Q_4\) is trace-class because
\(\sum_{n\ge1}|\chi_{-4}(n)|n^{-2}<\infty\). Its trace is therefore the
sum of the diagonal elements. The table \cref{eq:characters34-tables}
annuls the even pairs and alternates the odd ones:
\[
\Tr(N^{-2}Q_4)
=1-\frac1{3^2}+\frac1{5^2}-\frac1{7^2}+\cdots.
\]
The identity with \(L(2,\chi_{-4})\) is the definition of the
Dirichlet series; the equality with \(\operatorname{Im}\operatorname{Li}_2(\ii)\)
is obtained by taking the imaginary part of
\(\sum_{n\ge1}\ii^n/n^2\).
\end{proof}

The terminal value is
\[
G=0.9159655941772190150546035149\ldots.
\]
The determining structure is the identity \(J^2=-I\), which induces
exactly the alternating sign of the odd classes; the decimal
expansion is its terminal evaluation.

\section{Residual composition of the triadic--quaternary type}

The residual functional modulo three provides \(\chi_{-3}\). Its CRT
product with \(\chi_{-4}\) defines

\begin{equation}
\chi_{12}=\chi_{-3}\chi_{-4}.
\label{eq:chi12}
\end{equation}

On the units modulo twelve,

\[
\chi_{12}(1)=1,\quad
\chi_{12}(5)=-1,\quad
\chi_{12}(7)=-1,\quad
\chi_{12}(11)=1,
\]

and vanishes outside of them. Two exact evaluations are

\begin{align}
L(1,\chi_{12})&=\frac{\log(2+\sqrt3)}{\sqrt3},
\label{eq:L1chi12}\\
L(2,\chi_{12})&=\frac{\pi^2}{6\sqrt3}.
\label{eq:L2chi12}
\end{align}

The first contains the hyperbolic unit

\begin{equation}
\sqrt3L(1,\chi_{12})=\arcosh2=\log(2+\sqrt3).
\label{eq:pell2}
\end{equation}

Together with \(\log(3+2\sqrt2)=2\log(1+\sqrt2)\) and \(\log(30+\sqrt{899})\), this identity determines a
family of Pell scales whose members retain their respective algebraic
provenances.

Certified terminal evaluations are

\[
L(1,\chi_{12})=0.7603459963009463475310942549\ldots,
\]

\[
L(2,\chi_{12})=0.9497031262940093952634984917\ldots.
\]

\begin{proof}[Evaluations dodecaphasic]
Grouping by residue classes,
\begin{align*}
L(1,\chi_{12})
&=\sum_{n\ge0}\left(
\frac1{12n+1}-\frac1{12n+5}
-\frac1{12n+7}+\frac1{12n+11}\right)\\
&=\int_0^1
\frac{1-x^4-x^6+x^{10}}{1-x^{12}}\,dx
=\int_0^1\frac{1-x^4}{1+x^6}\,dx.
\end{align*}
A real decomposition into the quadratic factors of \(1+x^6\)
produces
\(\log(2+\sqrt3)/\sqrt3\). For the second moment, if
\(\psi_1\) is the trigamma function,
\begin{align*}
L(2,\chi_{12})
&=\frac1{144}\bigl[
\psi_1(1/12)-\psi_1(5/12)
-\psi_1(7/12)+\psi_1(11/12)\bigr]\\
&=\frac{\pi^2}{144}\left(
\csc^2\frac\pi{12}-\csc^2\frac{5\pi}{12}\right)
=\frac{\pi^2}{6\sqrt3},
\end{align*}
where has been used
\(\psi_1(z)+\psi_1(1-z)=\pi^2\csc^2(\pi z)\). The two proofs use
the CRT table and not a target decimal value.
\end{proof}

\section{Euler and Stieltjes Constants and the Theta--Mellin Transform}

Let
\[
\vartheta(t)=\sum_{n\in\Z}e^{-\pi n^2t},\qquad t>0.
\]
For \(\Re s>1\), the theta--Mellin functional recovers the convergent identity

\[
\pi^{-s/2}\Gamma\!\left(\frac{s}{2}\right)\zeta(s)
=\frac12\int_0^\infty
\bigl(\vartheta(t)-1\bigr)t^{s/2-1}\,dt.
\]

The meromorphic continuation is obtained by splitting the integral at \(t=1\), using
\(\vartheta(t)=t^{-1/2}\vartheta(1/t)\), and separating the polar terms. Its expansion at \(s=1\),

\begin{equation}
\zeta(s)=\frac1{s-1}+\sum_{j\ge0}
\frac{(-1)^j\gamma_j}{j!}(s-1)^j,
\label{eq:stieltjes}
\end{equation}

situates \(\gamma=\gamma_0\) and the Stieltjes constants in a single
family of renormalized moments. This organization expresses an
analytical common provenance, without attributing an interpretation that is physically identical to all its levels.

